% page 326 du fac-simile (image p-325 du scan)
% bloc 152.4 x 233.3 mm ; marge gauche 8.0 mm
\pgc{-12.700mm}{0.000mm}{
\l{\cel{5}318}
\l{}
\l{\sou{kronometro}. Horlojeto di qua la pendulo esas kompensita,}
\l{\cel{3}tale ke lu ne varias segun la temperaturo. - DEFIRS.}
\l{}
\l{kropo. (anat.) Sako o membrano quan la uceli havas sub}
\l{\cel{3}la guturo, ed en qua la nutrivi sejornas ante pasar}
\l{\cel{3}aden la stomako. - DE.}
\l{}
\l{kroso.- I. Bastono eklezio-pastorala, signo di la epis-}
\l{\cel{3}kopeso od abadeso, kun un extremajo kurvigita forme di}
\l{\cel{3}voluto, quan la abadi kustumis portar kun la voluto}
\l{\cel{3}adinterne, e la episkopi adextere. - II. Parto kurv-}
\l{\cel{3}igita per qua finas fusilo, pistolo, ye la parto opozita}
\l{\cel{3}a la tubo. - F.}
\l{}
\l{krotalo. (zool.)\cel{2}Serpento venenoza, di qua la korneti,}
\l{\cel{3}inter-inkastrita ye la extremajo di la kaudo, kliktas. -}
\l{\cel{3}L.\cel{1}\sou{crotalus}. - FISL.}
\l{}
\l{krotono. (bot.)\cel{2}Planto qua formacas genero di la familio}
\l{\cel{3}\textquotedbl{}euforbiacei\textquotedbl{}, di qua un speco donas la tornasolo, e}
\l{\cel{3}di qua un altra donas emetiko e purgivo, e grani de}
\l{\cel{3}qui onu ekpresas oleo akra, uzata kom drastiko e kom}
\l{\cel{3}revulsivo violenta. - L.\cel{1}\sou{croton.}\cel{2}- DEFIRS.}
\l{}
\l{krozar. (netrans.)\cel{2}(\sou{aludante navo}). Trairar maro, en ula}
\l{\cel{3}regioni, por surveyar li. - DEFIRS.}
\l{}
\l{kruco. I. (\sou{epoki antiqua}).(1) Sorto di jibeto, maxim-multa-}
\l{\cel{3}kaze trairata da traverso, sur qua onu igis mortar ula}
\l{\cel{3}krimininti, surligita o surklovagita ye la extremaji di}
\l{\cel{3}lia membri. (2) La jibeto sur qua Iesu Kristo esis klov-}
\l{\cel{3}agata ed ocidata (\sou{segun kristanismo}). (3).\cel{1}\sou{Metaf}.: La}
\l{\cel{3}religio di Iesu Kristo. (4) Objekto krucoforma, ek ligno,}
\l{\cel{3}metalo, e c., simbolo di Iesu Kristo krucagita. -}
\l{\cel{3}II. Distingivo honorizanta, ornivo ek kruco. - III.}
\l{\cel{3}To quo esas krucoforma. - DEFIRS.}
\l{}
\l{krucho. Vazo tera, terakota, larja-ventra, kun anso e kolo,}
\l{\cel{3}destinita a kontenar liquidi. - DFR.}
\l{}
\l{krucifera. (bot.)\cel{3}Di qua la petali esas krucoforma.}
\l{\cel{3}- DEFIRS.}
\l{}
\l{\sou{krucifixo}. Kruco ek ligno, metalo, ivoro, e c., sur qua}
\l{\cel{3}esas reprezentita Iesu Kristo krucagita. - DEFIRS.}
\l{}
\l{\sou{kru}da. I. (\sou{aludante la substanci nutrivala)}. Qua ne esas}
\l{\cel{3}koquita o bakita. - II. (\sou{aludante materio quan onu sub}-}
\l{\cel{3}\sou{misas a preparo)}.\cel{2}Qua ne subisis ta preparo. - II.}
\l{\cel{3}(metaf.) Quan nulo minfortigas (kom ex. : termino, dis-}
\l{\cel{3}kurso, kolori qui ne fuzesas kun olti qui cirkondas li).}
\l{\cel{3}- EFIS.}
}
